\section{Etika a regulace při sběru a uchovávání dat}

\subsection{Smysl etiky a regulace v datové analytice}
Datová analytika nepracuje jen s technickými záznamy. Data mohou popisovat lidi, jejich chování,
zdraví, finance, názory, polohu, práci nebo vztahy. Výstupy analýz mohou ovlivnit cenu služby,
schválení úvěru, výběr uchazeče o práci, cílení reklamy, interní kontrolu zaměstnanců nebo veřejné
rozhodování. Proto nestačí otázka, zda data umíme získat a zpracovat. Je nutné řešit také to, zda
je takové zpracování zákonné, přiměřené, bezpečné, vysvětlitelné a společensky přijatelné.

\term{Regulace} stanoví minimální závazná pravidla. V evropském prostředí jde hlavně o GDPR, český
zákon č. 110/2019 Sb., o zpracování osobních údajů, pravidla pro elektronickou komunikaci a nově také
AI Act u vybraných systémů umělé inteligence. \term{Etika} jde dál než samotná zákonnost. Ptá se,
zda je zpracování férové, zda lidé rozumějí dopadům, zda nedochází k diskriminaci, zda se s daty
nepracuje invazivně a zda analytik nepřekračuje účel, pro který byla data získána.

\subsection{Informační etika}
\term{Informační etika} je oblast morálky spojená se vznikem, šířením, transformací, ukládáním,
vyhledáváním, využíváním a organizací informací. Pro datovou analytiku je přímo použitelná: analytik
nepracuje s neutrálními čísly, ale s informacemi, které mohou mít dopad na lidi, organizace a veřejné
rozhodování.

Typické etické otázky:
\begin{itemize}
  \item je sběr dat přiměřený účelu, nebo se sbírá vše \uv{pro jistotu};
  \item vědí lidé, že jsou jejich údaje zpracovávány a k čemu;
  \item lze výstup analýzy rozumně vysvětlit člověku, kterého se týká;
  \item nevede použitý model k nepřímé diskriminaci určité skupiny;
  \item jsou data přesná a aktuální, nebo mohou poškodit člověka chybným závěrem;
  \item existuje odpovědná osoba, auditní stopa a možnost nápravy;
  \item je riziko úniku dat úměrně ošetřené.
\end{itemize}

\subsection{Etické teorie použitelné na data}
Etické dilema v datové analytice lze posuzovat z více pohledů. \term{Deontologický} přístup zdůrazňuje
práva a povinnosti, například respekt k souhlasu a k pravidlům zpracování bez ohledu na okamžitý
užitek. \term{Utilitarismus} hodnotí důsledky a ptá se, zda zpracování přinese co největší užitek
pro co nejvíce lidí. \term{Etika ctnosti} se ptá, zda analytik jedná poctivě, odpovědně a férově.
\term{Kontextuální přístup} připomíná, že stejný údaj může být přijatelný v jednom vztahu a
problematický v jiném, například zdravotní údaj u lékaře a v reklamním cílení.

U zkoušky je dobré ukázat, že právní soulad nestačí. Analýza může být formálně legální, ale pořád
eticky sporná, pokud je netransparentní, nepřiměřená, manipulační nebo vytváří nerovnováhu moci mezi
organizací a člověkem.

\subsection{Masonův rámec PAPA}
Pro základní strukturu odpovědi se hodí Masonův rámec PAPA: \textit{Privacy, Accuracy, Property,
Accessibility}.

\term{Privacy - soukromí} řeší, jaké informace si lidé mohou ponechat pro sebe, jaké informace se o
nich smí sbírat a za jakých podmínek. V analytice sem patří sledování chování uživatelů, spojování
dat z více zdrojů, práce s polohou, profilování a citlivé segmentace.

\term{Accuracy - přesnost} se týká správnosti, úplnosti a aktuálnosti informací. Chybná data mohou
vést k chybnému rozhodnutí. Typickým problémem je zastaralý záznam v CRM, špatné spárování osob,
duplicita, zkreslený trénovací dataset nebo model s vysokou chybovostí pro určitou skupinu.

\term{Property - vlastnictví a práva k informacím} řeší, kdo má právo data používat, sdílet, prodávat,
licencovat nebo zveřejnit. Patří sem autorská práva, databázová práva, obchodní tajemství, licence
datových sad, smluvní omezení API a práva subjektů údajů.

\term{Accessibility - dostupnost} řeší, kdo má k informacím přístup a kdo je z přístupu vyloučen.
V organizaci se promítá do řízení oprávnění, rolí, auditních logů, principu need-to-know a dostupnosti
dat pro oprávněné uživatele. Ve společnosti souvisí s digitální propastí a transparentností veřejných
informací.

\subsection{Mikroetické a makroetické problémy}
\term{Mikroetické problémy} se týkají chování jednotlivce nebo konkrétní organizace. Analytik by
neměl zneužívat přístup k datům, vyhledávat informace bez pracovního důvodu, manipulovat výsledky,
zamlčovat nejistotu, kopírovat data mimo schválené prostředí ani obcházet bezpečnostní pravidla.

\term{Makroetické problémy} se týkají širších dopadů dat a technologií na společnost. Patří sem
masové sledování, automatizované rozhodování bez možnosti obrany, diskriminace modelů, informační
asymetrie mezi firmou a zákazníkem, neprůhledné algoritmy, manipulativní personalizace a koncentrace
datové moci u několika velkých subjektů.

\subsection{Etika sběru dat}
Při sběru dat je potřeba řešit účel, přiměřenost, transparentnost a svobodnost rozhodnutí člověka.
Dobrý sběr dat má jasný důvod a předem stanovené hranice.

Zásady etického sběru:
\begin{itemize}
  \item sbírat jen data potřebná pro konkrétní účel;
  \item vysvětlit, kdo data sbírá, proč, jak dlouho je uchová a komu je předá;
  \item u dotazníků a rozhovorů respektovat dobrovolnost, anonymitu a citlivost odpovědí;
  \item nepoužívat nátlak, skryté účely, zavádějící otázky ani tmavé vzorce v rozhraní;
  \item oddělit identifikační údaje od analytických proměnných, pokud identifikace není nutná;
  \item nesbírat citlivé údaje jen proto, že mohou být \uv{někdy užitečné};
  \item zohlednit zranitelné skupiny, například děti, pacienty, zaměstnance nebo klienty veřejné správy.
\end{itemize}

Typickými riziky sběru jsou \term{sledování}, kdy lidé pod tlakem monitoringu mění chování, a
\term{dotazování} na citlivé věci, které může vyvolat nepohodlí nebo zkreslené odpovědi.
Ve výzkumu pomáhá transparentní plán: předem stanovené hypotézy, kritéria ukončení sběru a způsob
analýzy snižují riziko, že se hypotézy vytvoří až podle nalezených výsledků. V organizacích je potřeba
počítat i se sociálním inženýrstvím, například phishingem, pretextingem nebo vylákáním přístupu k
datům pod falešnou záminkou.

\subsection{Etika uchovávání a používání dat}
Při uchovávání dat je potřeba řešit životní cyklus dat: vytvoření, použití, sdílení, archivaci,
zálohování a smazání. Eticky i právně problematické je uchovávat osobní údaje neomezeně dlouho jen
proto, že úložiště je levné.

Důležité otázky:
\begin{itemize}
  \item kdo je za dataset odpovědný;
  \item kdo má přístup a zda je přístup pravidelně revidován;
  \item jak dlouho se data uchovávají a podle jakého pravidla se mažou;
  \item zda jsou data šifrovaná nebo pseudonymizovaná;
  \item zda existuje auditní stopa, kdo s daty pracoval;
  \item zda jsou zálohy chráněné stejně jako produkční data;
  \item zda výstupy neprozrazují osobní údaje nepřímo, například přes malé skupiny v agregaci.
\end{itemize}

\term{Kontextuální integrita} říká, že soukromí není jen skrytí informací, ale dodržení očekávaného
toku informací v dané situaci. Údaj sdělený bance kvůli posouzení úvěru nemá automaticky sloužit k
libovolnému marketingovému profilování. Podobně \term{řízení hranic soukromí} připomíná, že lidé
sdílejí informace za určitých podmínek a tyto hranice je potřeba při dalším použití respektovat.

\subsection{GDPR jako hlavní rámec ochrany osobních údajů}
\term{GDPR} je obecné nařízení EU o ochraně osobních údajů. Vztahuje se na zpracování osobních údajů
v souvislosti s činností správce nebo zpracovatele v EU a také na některé situace, kdy subjekt mimo
EU nabízí zboží či služby lidem v EU nebo sleduje jejich chování. V ČR na GDPR navazuje zákon
č. 110/2019 Sb., o zpracování osobních údajů. Dozorovým úřadem je Úřad pro ochranu osobních údajů
(\term{ÚOOÚ}).

\term{Osobní údaj} je jakákoli informace vztahující se k identifikované nebo identifikovatelné fyzické
osobě. Nemusí jít jen o jméno a rodné číslo. Osobním údajem může být e-mail, telefon, IP adresa,
cookie identifikátor, poloha, číslo zákazníka, fotografie, hlasový záznam nebo kombinace údajů, která
umožní člověka rozpoznat.

\term{Zpracování} je velmi široký pojem: sběr, záznam, uložení, strukturování, úprava, vyhledání,
použití, zpřístupnění, předání, kombinování, omezení, výmaz i zničení dat. Už samotné uložení osobních
údajů v databázi je zpracování.

\term{Subjekt údajů} je člověk, kterého se osobní údaje týkají. GDPR chrání fyzické osoby, ne právnické
osoby jako takové. Údaje o firmě nejsou samy o sobě osobní údaje, ale údaje o živnostníkovi, kontaktní
osobě, zaměstnanci nebo jednateli osobními údaji být mohou.

\subsection{Správce, zpracovatel a příjemce}
\term{Správce} určuje účely a prostředky zpracování. Rozhoduje tedy, proč a jak se osobní údaje
zpracovávají. Příkladem je e-shop, který zpracovává údaje zákazníků pro vyřízení objednávky, nebo
zaměstnavatel, který vede personální agendu.

\term{Zpracovatel} zpracovává osobní údaje pro správce a podle jeho pokynů. Příkladem může být
poskytovatel cloudového úložiště, externí mzdová firma, provozovatel e-mailingového nástroje nebo
dodavatel analytické platformy. Vztah správce a zpracovatele má být upraven smlouvou nebo jiným
právním aktem.

\term{Společní správci} společně určují účely a prostředky zpracování. Musí si transparentně rozdělit
odpovědnosti, zejména vůči subjektům údajů.

\term{Příjemce} je subjekt, kterému jsou osobní údaje zpřístupněny. Nemusí být vždy zpracovatelem.
Například dopravce, kterému e-shop předá adresu pro doručení, může být samostatným správcem pro vlastní
část činnosti, nikoli zpracovatelem e-shopu.

\subsection{Zvláštní kategorie osobních údajů}
GDPR přísněji chrání \term{zvláštní kategorie osobních údajů}. Patří sem údaje vypovídající o rasovém
nebo etnickém původu, politických názorech, náboženském nebo filozofickém přesvědčení, členství v
odborech, genetické údaje, biometrické údaje za účelem jedinečné identifikace, údaje o zdraví a údaje
o sexuálním životě nebo sexuální orientaci.

Zpracování těchto údajů je obecně zakázáno, pokud neplatí některá zvláštní výjimka, například výslovný
souhlas, pracovněprávní povinnost, ochrana životně důležitých zájmů, zdravotní péče, veřejné zdraví
nebo významný veřejný zájem stanovený právem. Pro analytiku to znamená, že citlivé proměnné nelze
jen tak přidat do modelu, i když by zvýšily predikční přesnost.

\subsection{Právní tituly zpracování}
Každé zpracování osobních údajů musí mít právní titul podle čl. 6 GDPR. Základní právní tituly jsou:
\begin{itemize}
  \item \term{souhlas} - svobodný, konkrétní, informovaný a jednoznačný projev vůle;
  \item \term{plnění smlouvy} - zpracování nutné pro smlouvu nebo kroky před jejím uzavřením;
  \item \term{právní povinnost} - například účetní, daňová nebo pracovněprávní evidence;
  \item \term{životně důležité zájmy} - ochrana života nebo zdraví;
  \item \term{veřejný zájem nebo výkon veřejné moci} - typicky veřejná správa;
  \item \term{oprávněný zájem} - zájem správce nebo třetí strany, pokud nepřevažují práva a svobody
  subjektu údajů.
\end{itemize}

Souhlas není univerzální řešení. Musí být odvolatelný a nesmí být vynucený. V pracovním vztahu bývá
problematický, protože zaměstnanec nemusí být vůči zaměstnavateli ve skutečně svobodném postavení.
U oprávněného zájmu je potřeba provést balanční test: existuje legitimní zájem, je zpracování nutné
a nepřevažují práva subjektů údajů?

\subsection{Principy GDPR}
GDPR stojí na několika principech. U zkoušky je dobré umět je nejen vyjmenovat, ale i vysvětlit na
příkladu.

\begin{center}
\small
\begin{tabular}{@{}>{\raggedright\arraybackslash}p{0.28\textwidth}
                >{\raggedright\arraybackslash}p{0.61\textwidth}@{}}
\hline
\term{Princip} & \term{Význam v analytice} \\
\hline
Zákonnost, korektnost a transparentnost & zpracování má právní titul, není klamavé a člověk rozumí,
co se s jeho údaji děje \\
Účelové omezení & data získaná pro jeden účel nelze automaticky použít pro jiný neslučitelný účel \\
Minimalizace údajů & sbírají se jen údaje potřebné pro danou úlohu, ne maximum dostupných atributů \\
Přesnost & nepřesné nebo zastaralé údaje se opravují nebo mažou \\
Omezení uložení & osobní údaje se neuchovávají déle, než je potřeba \\
Integrita a důvěrnost & data jsou chráněna proti neoprávněnému přístupu, ztrátě, zničení a změně \\
Odpovědnost & správce musí být schopen prokázat soulad s pravidly, nestačí tvrdit, že je dodržuje \\
\hline
\end{tabular}
\end{center}

\subsection{Informační povinnost a práva subjektů údajů}
Transparentnost znamená, že subjekt údajů má dostat srozumitelné informace o zpracování. Typicky jde
o identitu správce, účely zpracování, právní titul, kategorie údajů, příjemce, dobu uložení, práva
subjektu údajů, možnost stížnosti u dozorového úřadu a informaci o automatizovaném rozhodování nebo
profilování, pokud se používá.

Základní práva subjektu údajů:
\begin{itemize}
  \item právo na informace;
  \item právo na přístup k údajům;
  \item právo na opravu;
  \item právo na výmaz, někdy označované jako \uv{právo být zapomenut};
  \item právo na omezení zpracování;
  \item právo na přenositelnost údajů;
  \item právo vznést námitku;
  \item právo nebýt za určitých podmínek předmětem rozhodnutí založeného výhradně na automatizovaném
  zpracování, včetně profilování, pokud má právní nebo obdobně významné účinky.
\end{itemize}

V analytickém projektu je proto nutné vědět, odkud data pocházejí, jaký je jejich účel, jak vyřídit
žádost o přístup nebo výmaz a zda by výmaz jednoho člověka nerozbil auditovatelnost modelu. Tyto
otázky se mají řešit při návrhu, ne až po obdržení žádosti.

\subsection{Ochrana dat od návrhu a ve výchozím nastavení}
\term{Privacy by design} znamená, že ochrana osobních údajů je součástí návrhu systému od začátku.
Nejde o dodatečné právní prohlášení po dokončení databáze. Příkladem je oddělení identifikačních
údajů od analytických atributů, pseudonymizace, omezení přístupů, auditní logy, šifrování a pravidla
retence.

\term{Privacy by default} znamená, že výchozí nastavení má být co nejšetrnější. Systém nemá automaticky
zveřejňovat profil, sdílet data se všemi odděleními, uchovávat údaje navždy nebo zapínat marketingové
sledování bez vhodného právního titulu.

\subsection{Anonymizace a pseudonymizace}
\term{Pseudonymizace} znamená, že se přímé identifikátory nahradí kódem nebo aliasem a dodatečné
informace potřebné k přiřazení zpět ke konkrétnímu člověku se uchovávají odděleně. Pseudonymizovaná
data jsou stále osobní údaje, protože při použití dodatečných informací je možné osobu znovu určit.

\term{Anonymizace} znamená, že jednotlivce nelze identifikovat žádným rozumně pravděpodobným prostředkem.
Pokud je anonymizace provedena správně, GDPR se na anonymizovaná data už nevztahuje. V praxi je ale
anonymizace obtížná: kombinace věku, pohlaví, obce, zaměstnání a vzácné diagnózy může člověka odhalit
i bez jména.

Praktické techniky:
\begin{itemize}
  \item odstranění přímých identifikátorů;
  \item nahrazení ID pseudonymem nebo hashem se solí;
  \item generalizace hodnot, například věk na věkové skupiny;
  \item agregace do skupin s dostatečným počtem osob;
  \item potlačení vzácných kategorií;
  \item přidání šumu nebo použití metod typu differential privacy u citlivých agregací.
\end{itemize}

\subsection{Bezpečnost a uchovávání osobních údajů}
GDPR vyžaduje bezpečnost přiměřenou riziku. Neexistuje jeden univerzální seznam opatření pro všechny
systémy. Jiná opatření stačí pro newsletter s e-maily a jiná pro zdravotní data, biometrická data
nebo úvěrové skórování.

Typická technická a organizační opatření:
\begin{itemize}
  \item řízení přístupů podle rolí a princip nejmenších oprávnění;
  \item vícefaktorové ověřování u citlivých systémů;
  \item šifrování dat při přenosu a podle rizika i při uložení;
  \item pseudonymizace analytických datasetů;
  \item oddělení produkčních, testovacích a analytických prostředí;
  \item logování a pravidelná kontrola přístupů;
  \item zálohy, test obnovy a plán reakce na incident;
  \item smluvní a bezpečnostní kontrola dodavatelů;
  \item školení lidí, kteří s osobními údaji pracují.
\end{itemize}

\term{Doba uložení} má být odvozena od účelu a právních povinností. Účetní doklady se uchovávají
podle zákonných lhůt, zákaznické chování pro personalizaci jen tak dlouho, jak je pro účel nezbytné,
a testovací kopie databází by neměly obsahovat reálné osobní údaje, pokud to není nutné.

\subsection{Porušení zabezpečení osobních údajů}
\term{Porušení zabezpečení osobních údajů} je incident vedoucí k náhodnému nebo protiprávnímu zničení,
ztrátě, změně, neoprávněnému zpřístupnění nebo přístupu k osobním údajům. Může jít o hackerský útok,
ztracený notebook, omylem odeslaný e-mail, špatně nastavené cloudové úložiště nebo interní neoprávněný
přístup.

Správce má incident posoudit podle rizika pro práva a svobody lidí. Pokud je pravděpodobné riziko,
ohlašuje incident dozorovému úřadu bez zbytečného odkladu, pokud možno do 72 hodin od zjištění.
Pokud je pravděpodobné vysoké riziko pro dotčené osoby, musí být informovány i ony. Zpracovatel
oznamuje incident správci bez zbytečného odkladu.

\subsection{DPIA a pověřenec pro ochranu osobních údajů}
\term{DPIA - Data Protection Impact Assessment} neboli posouzení vlivu na ochranu osobních údajů je
analýza rizik pro práva a svobody lidí. Provádí se před zahájením zpracování, pokud zpracování
pravděpodobně povede k vysokému riziku. Typicky půjde o rozsáhlé zpracování citlivých údajů, rozsáhlé
systematické monitorování nebo automatizované hodnocení osob s významnými dopady.

DPIA má popsat:
\begin{itemize}
  \item účel a rozsah zpracování;
  \item nezbytnost a přiměřenost zpracování;
  \item rizika pro subjekty údajů;
  \item technická a organizační opatření ke snížení rizik;
  \item zbytkové riziko a odpovědnost za rozhodnutí.
\end{itemize}

\term{Pověřenec pro ochranu osobních údajů} neboli DPO je povinný například u orgánů veřejné moci,
u rozsáhlého pravidelného a systematického monitorování osob nebo u rozsáhlého zpracování zvláštních
kategorií údajů. DPO radí správci nebo zpracovateli, monitoruje soulad s pravidly a spolupracuje
s dozorovým úřadem.

\subsection{Předávání osobních údajů mimo EU/EHP}
Předání osobních údajů do třetí země nebo mezinárodní organizaci je možné jen tehdy, pokud není
oslabena úroveň ochrany zaručená GDPR. Prakticky se používají hlavně:
\begin{itemize}
  \item \term{rozhodnutí o odpovídající ochraně} vydané Evropskou komisí;
  \item \term{standardní smluvní doložky} a posouzení rizik konkrétního předání;
  \item \term{závazná podniková pravidla} pro skupiny podniků;
  \item omezené výjimky pro specifické situace, například výslovný souhlas nebo nutnost pro smlouvu.
\end{itemize}

Pro analytickou praxi je důležité, že použití cloudové služby nebo externího API může znamenat předání
dat mimo EU/EHP, i když uživatel sedí v ČR. Nestačí se dívat jen na sídlo dodavatele; je potřeba vědět,
kde se data ukládají, kdo k nim má přístup, zda jsou zapojeni subdodavatelé a jaké jsou smluvní záruky.

\subsection{Elektronická komunikace, cookies a marketing}
Vedle GDPR je u online sběru dat důležitá také regulace elektronické komunikace. Cookies, reklamní
identifikátory, newslettery a behaviorální reklama často kombinují pravidla ochrany osobních údajů,
pravidla pro elektronické komunikace a ochranu spotřebitele.

Prakticky:
\begin{itemize}
  \item technicky nezbytné cookies lze obvykle použít bez souhlasu;
  \item analytické, reklamní a profilovací cookies typicky vyžadují předchozí souhlas;
  \item souhlas má být svobodný, informovaný, konkrétní a odvolatelný;
  \item obchodní sdělení musí respektovat pravidla pro nevyžádanou komunikaci;
  \item u profilování pro marketing je nutné počítat s právem vznést námitku.
\end{itemize}

\subsection{Etika algoritmů a modelů}
U datové analytiky a strojového učení vznikají zvláštní rizika. Model může být formálně přesný v
průměru, ale nespravedlivý pro určitou skupinu. Může používat proxy proměnné, které nepřímo nahrazují
chráněnou charakteristiku, například adresu jako zástupný znak pro socioekonomický status nebo etnické
složení lokality.

Důležité principy:
\begin{itemize}
  \item \term{fairness} - kontrolovat, zda model neznevýhodňuje určité skupiny;
  \item \term{explainability} - umět vysvětlit podstatné důvody rozhodnutí nebo alespoň chování modelu;
  \item \term{accountability} - mít odpovědnou osobu, dokumentaci, audit a možnost nápravy;
  \item \term{human oversight} - u významných rozhodnutí ponechat smysluplný lidský dohled;
  \item \term{robustness} - sledovat chyby, drift, zneužití a bezpečnost modelu;
  \item \term{contestability} - umožnit člověku rozhodnutí napadnout nebo požádat o přezkum.
\end{itemize}

GDPR obsahuje pravidla pro automatizované individuální rozhodování včetně profilování. Subjekt údajů
má za určitých podmínek právo nebýt předmětem rozhodnutí založeného výhradně na automatizovaném
zpracování, pokud má právní nebo obdobně významné účinky. Pokud se takové rozhodování použije na
základě výjimky, musí existovat vhodná ochranná opatření, například možnost lidského zásahu, vyjádření
stanoviska a napadení rozhodnutí.

AI Act doplňuje pro vysoce rizikové systémy AI další požadavky, například řízení rizik, kvalitu dat,
dokumentaci, logování, transparentnost vůči nasazovateli, lidský dohled, přesnost, robustnost a
kybernetickou bezpečnost. U datové analytiky je podstatné hlavně to, že kvalita trénovacích,
validačních a testovacích dat, kontrola biasu a dokumentace nejsou jen technické detaily, ale součást
řízení rizik.

\subsection{Praktický příklad}
Firma chce vytvořit model pro predikci odchodu zákazníků. Má interní transakce, komunikaci se zákazníky,
reklamace a webové chování.

Správný postup:
\begin{enumerate}
  \item určit účel modelu, například retenční nabídka stávajícím zákazníkům;
  \item ověřit právní titul pro jednotlivé zdroje dat;
  \item nepoužívat citlivé údaje, pokud pro to není jasný právní důvod;
  \item minimalizovat atributy a omezit dobu historie;
  \item pseudonymizovat analytický dataset;
  \item zkontrolovat, zda model nepoužívá problematické proxy proměnné;
  \item dokumentovat zdroje dat, transformace, metriky modelu a omezení;
  \item nastavit přístupová práva, retenci, monitoring driftu a proces pro žádosti subjektů údajů;
  \item vysvětlit v interních pravidlech, kdo smí výstup použít a k jakému rozhodnutí.
\end{enumerate}

Chybný postup by byl stáhnout maximum dostupných dat, přidat je do modelu bez jasného účelu, výsledek
použít pro agresivní marketing a údaje uchovávat bez stanovené doby výmazu.

\pozor
Nejčastější chyba u zkoušky je tvrdit, že GDPR znamená \uv{všechno jen se souhlasem}. Správně je:
každé zpracování potřebuje právní titul, souhlas je jen jeden z nich a často není nejvhodnější. Druhá
častá chyba je zaměnit pseudonymizaci za anonymizaci. Pseudonymizovaná data jsou stále osobní údaje.

\subsection{Shrnutí ke zkoušce}
\begin{itemize}
  \item Etika řeší férovost, přiměřenost, dopady a odpovědnost; regulace stanoví závazná pravidla.
  \item Masonův rámec PAPA pokrývá soukromí, přesnost, vlastnictví a dostupnost informací.
  \item GDPR chrání osobní údaje fyzických osob a vztahuje se na široce chápané zpracování.
  \item Správce určuje účel a prostředky zpracování, zpracovatel jedná pro správce podle jeho pokynů.
  \item Každé zpracování musí mít právní titul a respektovat principy GDPR: zákonnost, transparentnost,
  účelové omezení, minimalizaci, přesnost, omezení uložení, integritu, důvěrnost a odpovědnost.
  \item Zvláštní kategorie údajů, například zdravotní nebo biometrické údaje, vyžadují zvláštní opatrnost.
  \item Subjekty údajů mají práva na informace, přístup, opravu, výmaz, omezení, přenositelnost,
  námitku a ochranu před některým automatizovaným rozhodováním.
  \item Bezpečnost zahrnuje technická i organizační opatření, retenci, řízení přístupů, šifrování,
  audit a reakci na incidenty.
  \item DPIA se provádí u rizikového zpracování, zejména u rozsáhlého monitoringu, citlivých údajů
  nebo významného automatizovaného hodnocení osob.
  \item Anonymizace odstraňuje identifikovatelnost, pseudonymizace ji jen omezuje.
  \item U algoritmů je nutné řešit bias, vysvětlitelnost, lidský dohled, auditovatelnost a možnost nápravy.
\end{itemize}

\subsection{Zdroje}
\begin{itemize}
  \item EUR-Lex: \href{https://eur-lex.europa.eu/legal-content/CS/TXT/?uri=CELEX:32016R0679}{Nařízení (EU) 2016/679 - GDPR}.
  \item Evropská komise: \href{https://commission.europa.eu/law/law-topic/data-protection/rules-business-and-organisations/obligations/controllerprocessor/what-data-controller-or-data-processor_en}{What is a data controller or a data processor?}.
  \item European Data Protection Board: \href{https://www.edpb.europa.eu/sme-data-protection-guide/faq-frequently-asked-questions/answer/what-difference-between_en}{Pseudonymised data and anonymised data}.
  \item ÚOOÚ: \href{https://uoou.gov.cz/cinnost/ochrana-osobnich-udaju}{Ochrana osobních údajů}.
  \item ÚOOÚ: \href{https://uoou.gov.cz/profesional/qa-otazky-a-odpovedi/posouzeni-vlivu-na-ochranu-osobnich-udaju}{DPIA - Posouzení vlivu na ochranu osobních údajů}.
  \item ÚOOÚ: \href{https://uoou.gov.cz/profesional/poruseni-zabezpeceni-osobnich-udaju}{Porušení zabezpečení osobních údajů}.
  \item EUR-Lex: \href{https://eur-lex.europa.eu/legal-content/CS/TXT/?uri=CELEX:32024R1689}{Nařízení (EU) 2024/1689 - AI Act}.
  \item OECD: \href{https://www.oecd.org/en/topics/ai-principles.html}{AI Principles}.
\end{itemize}

\newpage
